\section{Exact height: termination and complexity}
We use division-free arithmetic circuits over $\mathbb F_p$, with the
$d$ variables explicitly listed, for a homogeneous degree-$d$ equation
of a promised smooth hypersurface $X\subset\mathbb P^{d-1}$, $d\ge3$.
The output of the exact-height problem is a binary integer or $\infty$.
The prime is fixed in the modular counting assertions; the polynomial-space
assertion also permits a binary-encoded input prime.  Set $n=d-2$.

\begin{proposition}[A finite stopping bound]
Every finite height of $X$ is at most
\begin{equation}
 B_d=\begin{cases}
 2&d=3,\\
 36&d=4,\\
 \binom{2d-1}{d-1}&d\ge5.
 \end{cases}
\label{eq:HC1}
\end{equation}
\end{proposition}
\begin{proof}
Van der Geer and Katsura (2003, Corollary~2.3, p.~100) prove
$h\le h^{n-1}(\Omega_X^1)+1$ for finite $h$.  Serre duality and
$\omega_X\simeq\mathcal O_X$ identify this dimension with $h^1(T_X)$.
For $n\ge3$, the restricted Euler sequence and hypersurface line-bundle
cohomology give $H^1(T_{\mathbb P^{d-1}}|_X)=0$: indeed
$H^1(\mathcal O_X(1))=H^2(\mathcal O_X)=0$.  The normal sequence
then gives
\[
 h^1(T_X)\le h^0(\mathcal O_X(d))
 =\binom{2d-1}{d-1}-1.
\]
For $d=4$ the same Euler sequence gives
$h^1(T_{\mathbb P^3}|_X)=h^2(\mathcal O_X)=1$, so that
$h^1(T_X)\le34+1=35$.  For $d=3$,
$h^0(\Omega_X^1)=1$.  This proves \eqref{eq:HC1}, without assuming that
infinitesimal automorphisms vanish or that $p$ does not divide $d$.
\end{proof}

\begin{proposition}[Exact modular coefficient algorithm]
Let $\widetilde F$ be any homogeneous integral lift and put
\[
 a_{p^j}=[(x_1\cdots x_d)^{p^j-1}]\widetilde F^{p^j-1}.
\]
Compute $a_{p^j}\bmod p^j$ successively for $1\le j\le B_d$.
The first nonzero residue has index $h(X)$; if all residues vanish,
$h(X)=\infty$.
\end{proposition}
\begin{proof}
Lemma~\ref{lem:logarithm-height}
says that conditional on $h\ge j$ one has
$v_p(a_{p^j})\ge j-1$, with equality exactly when $h=j$.
Thus the first nonzero residue has precisely the claimed meaning.
If all residues vanish, a finite height would exceed \eqref{eq:HC1}, which is
impossible.  Only precision $p^j$ is required at the $j$th step.

For a circuit implementation, first form the homogeneous components
through degree $d$ and keep the degree-$d$ output. Choose integer
representatives of its scalar literals once and fix this homogeneous
integral circuit lift for all steps. A literal lift of the original
circuit need not be homogeneous, because cancellations in unwanted
degrees may have occurred only modulo $p$.
Evaluate the resulting circuit in the finite ring
\[
 (\mathbb Z/p^j\mathbb Z)[x_1,\ldots,x_d]/
 (x_1^{p^j},\ldots,x_d^{p^j}),
\]
raise its value to $p^j-1$ by repeated squaring, and read the diagonal
coefficient. Nonnegative exponents make truncation exact.
\end{proof}

\subsection{The quasi-\texorpdfstring{$F$}{F}-splitting matrix and complexity}
\label{subsec:quasi-height-complexity}

The full-height iteration used here is an existing algorithm.
Kawakami, Takamatsu, and Yoshikawa give its exact criterion
\citep[Theorem~C, version~3, p.~4; Theorem~5.8, pp.~37--39]{KTY2025};
the criterion was already present in version~2 of their preprint, p.~5.
Batubara, Garzella, and Pan identify the fixed finite-dimensional operator
and implement its matrix iteration
\citep[Section~3.4 and Algorithm~4, p.~7]{BGP2025}.
Yobuko identifies quasi-$F$-split height with Artin--Mazur height for
Calabi--Yau varieties
\citep[Theorem~4.5; author version, pp.~9--10]{Yobuko2019}.
We analyze this algorithm under the circuit input convention of this paper.

Put $S=\mathbb F_p[x_1,\ldots,x_d]$ and write
\[
 D=d(p-1),\qquad g=F^{p-1},\qquad
 V=S_D,\qquad a=(p-1,\ldots,p-1),
\]
\[
 W=\dim_{\mathbb F_p}V=\binom{dp-1}{d-1}.
\]
Thus $V=S_D$ is the space of homogeneous degree-$D$ polynomials.
Choose a homogeneous integral coefficient lift
$G=\sum_{|\beta|=D}\widetilde c_\beta x^\beta$ of $g$, and define
\begin{equation}\label{eq:quasi-delta}
 \Delta=\left(
 \frac{G^p-\sum_{|\beta|=D}\widetilde c_\beta^{\,p}x^{p\beta}}p
 \right)\bmod p.
\end{equation}
This is the monomial representative $\Delta_1(g)$ of the cited criterion;
see also \citet[Proposition~3.2, Remark~3.4 and Algorithm~1, p.~5]{BGP2025}.
The sum is over distinct collected monomials.
The numerator is divisible by $p$ coefficientwise. Moreover, replacing
$G$ by $G+pH$, or a scalar lift by $\widetilde c_\beta+pb$, changes the
corresponding $p$th power by a multiple of $p^2$. Thus
\eqref{eq:quasi-delta} is independent of the chosen integer coefficient
lifts. This includes $p=2$ and requires no Teichm\"uller representatives.
Computationally, form the numerator modulo $p^2$, choose each residue in
$[0,p^2)$, and divide it by $p$; the resulting residue in $[0,p)$ is exact.

Define the $\mathbb F_p$-linear map $U$ on monomials by
\[
 U(x^\gamma)=
 \begin{cases}
 x^{(\gamma-a)/p},&\gamma\equiv a\pmod p,\\
 0,&\text{otherwise}.
 \end{cases}
\]
Since $\deg\Delta=pD$, the map
\[
 M:V\longrightarrow V,\qquad v\longmapsto U(\Delta v)
\]
preserves $V$: its retained monomials have degree
$(pD+D-d(p-1))/p=D$. In the monomial basis,
\begin{equation}\label{eq:quasi-matrix-entry}
 M_{\alpha,\beta}=[x^{p\alpha+a-\beta}]\Delta,
 \qquad |\alpha|=|\beta|=D,
\end{equation}
where a negative exponent gives zero. Let $\ell(v)=[x^a]v$.
The only degree-$D$ monomial outside $(x_1^p,\ldots,x_d^p)$ is $x^a$.
Consequently the cited height criterion gives
\begin{equation}\label{eq:quasi-height-iteration}
 h(X)=\inf\{j\ge1:\ell(M^{j-1}g)\ne0\},\qquad \inf\varnothing=\infty.
\end{equation}

For the field hypotheses, apply Kawakami, Takamatsu, and Yoshikawa's
projective criterion and Yobuko's
comparison after extension to $\overline{\mathbb F}_p$. The Cartier map
over a perfect field takes $p$th roots of coefficients, but every
coefficient of this recurrence lies in $\mathbb F_p$, so the recurrence
is exactly the linear one displayed above. Formal-group height is
unchanged by extending the base field, since nonvanishing of the first
nonzero coefficient of the $p$-series is unchanged. The adjunction and
cohomology calculations already made verify the Calabi--Yau hypotheses.
This applies also when $d=3$.

\begin{lemma}[A matrix stopping bound]
\label{lem:quasi-matrix-stopping}
In \eqref{eq:quasi-height-iteration}, a finite height is at most $W$.
Thus testing $j=1,\ldots,W$ detects infinite height as well. For the
smooth inputs here one may use the smaller cutoff $\min(W,B_d)$.
\end{lemma}
\begin{proof}
The sequence $y_r=\ell(M^rg)$ satisfies a linear recurrence of order
at most $W$ by Cayley--Hamilton. If its first $W$ terms vanish, the
recurrence forces every subsequent term to vanish. Otherwise the first
nonzero term has index at most $W-1$, which proves the assertion by
\eqref{eq:quasi-height-iteration}. This argument requires no computation
of the characteristic polynomial. The alternative bound $B_d$ is
\eqref{eq:HC1}.
\end{proof}

Write $\mathrm{FPSPACE}$ for function problems computable deterministically
with a number of work bits polynomial in the input length.

\begin{theorem}[Resource bounds for exact height]
\label{thm:quasi-height-resource-bounds}
For a binary-encoded prime $p$ and a division-free circuit of size $s$
in $d\ge3$ explicitly listed variables, promised to define a smooth
degree-$d$ hypersurface in $\mathbb P^{d-1}$, exact Artin--Mazur height,
including $\infty$, belongs to $\mathrm{FPSPACE}$.
There is also a deterministic dense implementation with bit complexity
\[
 2^{O(d\log p)}\operatorname{poly}(s,d,\log p).
\]
The polynomial-space and dense-time bounds refer to separate
implementations of the same height criterion.
\end{theorem}
\begin{proof}
Let $L$ denote input length and set $E=pD=dp(p-1)$. All scalar
calculations below take place modulo $p$ or $p^2$, so each scalar has
$O(\log p)$ bits. A monomial of total degree at most $E$ has an exponent
vector using $O(d\log(E+1))$ bits. Furthermore,
\[
 \log W=O(d(1+\log p)),\qquad
 \log(E+1)=O(\log d+\log p).
\]
These lengths are polynomial in $L$, because $d$ is explicitly bounded
by input length.

\paragraph{Coefficient access.}
For a degree-$d$ output coefficient of $F$, recurse through its arithmetic
circuit. At a product gate, enumerate all coordinatewise splits of the
target exponent vector, obtain the two child coefficients, and accumulate
their product modulo $p$. At a sum gate, add the child coefficients.
Constants and variables are immediate. A recursion path has at most $s$
gates, and each frame stores polynomial-length exponent vectors, gate
addresses, a loop counter and an $O(\log p)$-bit accumulator. Intermediate
gates need not be homogeneous. Equivalently, this computation evaluates
the circuit modulo all monomials of total degree greater than $d$, which
retains the required output exactly. The degree-$d$ output is selected
using the promise; an output coefficient of any other total degree is zero.

Coefficients of a power $A^r$ are obtained by the same convolution
recursion applied to
\[
 A^r=A^{\lfloor r/2\rfloor}A^{\lceil r/2\rceil},
\]
with base cases $r=0,1$. The power recursion has depth
$O(\log(r+1))$. First compute coefficients of $g=F^{p-1}$ modulo $p$.
Treat their representatives in $[0,p)$ as the coefficients of a virtual
homogeneous lift $G$, and compute coefficients of $G^p$ modulo $p^2$.
The two powering recursions add only $O(\log p)$ depth to the original
circuit recursion. In the pure-power sum in \eqref{eq:quasi-delta}, a
requested exponent has at most one contributing term, determined by
coordinatewise divisibility by $p$. Scalar powering modulo $p^2$ and
the exact scalar division by $p$ complete the coefficient calculation.
Thus every coefficient of $\Delta$, and hence every matrix entry in
\eqref{eq:quasi-matrix-entry}, is accessible in polynomial space.

\paragraph{Matrix powers.}
For $r\ge2$, compute one entry recursively from
\[
 (M^r)_{\alpha,\beta}
   =\sum_{|\gamma|=D}
      (M^{\lfloor r/2\rfloor})_{\alpha,\gamma}
      (M^{\lceil r/2\rceil})_{\gamma,\beta}\pmod p.
\]
The cases $r=0,1$ use the identity and the entry oracle just constructed.
For $r<W$, the matrix recursion has depth $O(\log W)$.
Every frame stores two index vectors, the current summation index, a
binary exponent, and one field accumulator. All have polynomial length.
An entry-oracle call adds only the polynomial stack used for coefficients.
Finally, compute
\[
 \ell(M^{j-1}g)=
 \sum_{|\beta|=D}(M^{j-1})_{a,\beta}g_\beta\pmod p
\]
and scan $j=1,\ldots,W$ using a binary counter. Return the first nonzero
test, or $\infty$ if all vanish. Lemma~\ref{lem:quasi-matrix-stopping}
proves correctness. The outer scan does not create a recursion chain of
length $W$; every requested power uses divide and conquer. The complete
stack, counters, exponents and scalar accumulators use $L^{O(1)}$ space.

\paragraph{A separate dense time bound.}
The number of monomials of total degree at most $E$ is
\[
 V_E=\binom{E+d}{d}
 \le \bigl(e[p(p-1)+1]\bigr)^d
 \le (ep^2)^d,
 \qquad W\le\binom{dp}{d}\le(ep)^d.
\]
Expand the original circuit modulo $p$ through total degree $d$,
using at most $\binom{2d}{d}\le4^d$ positions. Dense binary powering
computes $g$ and $G^p$ in $O(\log p)$ polynomial-product stages.
Naive coefficient convolution uses at most $V_E^2$ scalar products
per stage. Form $\Delta$ by \eqref{eq:quasi-delta} and construct the
$W^2$ matrix entries by coefficient lookup. At most $W-1$ matrix-vector
products then test the whole sequence, using $O(W^3)$ field operations.
Degree indexing and scalar arithmetic contribute a fixed polynomial in
$d,\log p$; even sequential array-access implementations contribute
only a fixed polynomial in $V_E,W$. Thus the total bit complexity is
a fixed polynomial in $s,d,\log p,V_E,W$, giving
$2^{O(d\log p)}\operatorname{poly}(s,d,\log p)$ uniformly in all
parameters. This dense implementation uses exponentially large arrays,
whereas the preceding recursion recomputes entries to obtain the
polynomial-space bound.
\end{proof}



\subsection{Uniform characteristic and exact counting}
Varying the characteristic gives a stronger lower bound than fixing one modular counting class. We state the encoding explicitly, since this distinction is essential.

The symbols $\mathrm P$ and $\mathrm{FP}$ denote deterministic
polynomial-time decision and function computation, respectively.
A superscript denotes access to the indicated oracle; in particular,
a $\#\mathrm P$ oracle returns an exact integer count.

\begin{theorem}[Uniform height two and exact counting]\label{thm:uniform-counting}
Let $\mathcal H_2$ be the promise decision problem whose input is a prime $p$ in unary and a width-three formula $\Phi$, promised to have $p\mid\#\operatorname{SAT}(\Phi)$, and whose answer is whether the counting hypersurface $X_\Phi$ has height two. Then
\[
 \mathrm{FP}^{\mathcal H_2}=\mathrm{FP}^{\#\mathrm P},
\]
where oracle calls to the promise problem must satisfy its promise.
The reduction from exact counting uses polynomially many nonadaptive calls. Consequently the full exact-height function for the circuit presentations of this paper is $\#\mathrm P$-hard under polynomial-time nonadaptive oracle reductions, even when every query has height at least two and its characteristic is polynomially bounded in the size of the original counting input.
\end{theorem}
\begin{proof}
For the upper inclusion, the ordered local lists defining $T_{p+1}(\Phi)$ have length polynomial in the unary input length and are checked by the displayed marginal equations. Hence this integer is a $\#\mathrm P$ count uniformly in $p$. On the promise, the oracle answer is obtained from its exact value by reducing modulo $p^2$ and testing whether the residue is nonzero. The height-two criterion proves correctness. An exact-counting oracle therefore simulates each promised $\mathcal H_2$ call in polynomial time.

For the other inclusion let $\Phi$ have $N\ge1$ explicitly declared variables, and put $S=\#\operatorname{SAT}(\Phi)$, so $0\le S\le2^N$. For any polynomially bounded positive prime $p$ and any $a\in\{0,\ldots,p-1\}$, construct a width-three formula $\Phi_a$ with exactly $S+a$ satisfying assignments as follows. Introduce a branch bit $b$ and $k=\lceil\log_2 p\rceil$ bits $z$, and use the Boolean predicate
\[
 (\neg b\wedge\Phi(x)\wedge z=0)
 \ \vee\
 (b\wedge x=0\wedge z<a).
\]
The two branches have respectively $S$ and $a$ assignments. The comparison $z<a$ has a Boolean circuit of size polynomial in $k$; the whole predicate therefore has polynomial size. Give each circuit gate a new variable and impose its truth table by clauses of width at most three, then force the output to one. For example, $v=u\wedge w$ is enforced by
\[
 (\neg v\vee u)\wedge(\neg v\vee w)
       \wedge(v\vee\neg u\vee\neg w),
\]
and a negation uses two binary clauses. OR gates are treated by the dual three clauses. Each input assignment has exactly one extension to the gate variables, so this transformation preserves the satisfying count exactly. The case $a=0$ uses the always-false second branch. Declared variables and auxiliary variables remain part of the counted assignment.

For each prime use the height-two gadget $\Psi_p$ already proved above: for odd $p$ take $\Theta_{p-1,p-2}$, and for $p=2$ take the displayed three-clause characteristic-two gadget. It has $O(p^2)$ clauses, $p\mid\#\operatorname{SAT}(\Psi_p)$, and
\[
 v_p\bigl(T_{p+1}(\Psi_p)\bigr)=1.
\]
Conjoin $\Phi_a$ with $\Psi_p$ on disjoint variable sets. Every resulting query satisfies the Hasse-zero promise. Multiplicativity and $T_{p+1}(\Phi_a)\equiv(S+a)^2\pmod p$ give
\[
 h(X_{\Phi_a\wedge\Psi_p})=2
 \quad\Longleftrightarrow\quad S+a\not\equiv0\pmod p.
\]
Exactly one of the $p$ answers is therefore negative, identifying $S\bmod p$. These $p$ queries are fixed before any answer is read.

We give an elementary polynomial bound for the primes required. Put $M=4(N+1)^2$, let $P_M$ be the product of the primes at most $M$, and let $L_M=\operatorname{lcm}(1,\ldots,M)$. Every binomial coefficient $\binom Mj$ divides $L_M$: for each prime, its factorial-valuation formula is a sum of at most $\lfloor\log_p M\rfloor$ terms, each zero or one. Hence
\[
 L_M\ge\max_j\binom Mj\ge\frac{2^M}{M+1},
 \qquad L_M\le P_M^{\lfloor\log_2M\rfloor}.
\]
For $N\ge1$, $M\ge16$, $\log_2(M+1)\le M/2$, and $\log_2M\le2(N+1)$. These inequalities imply $P_M\ge2^{N+1}>2^N$. Trial division or a sieve lists all primes up to $M$ in polynomial time. Asking all shifts for all these primes takes at most $M^2$ queries. Chinese remaindering then recovers the unique $S\in[0,2^N]$ from its residues. The case $N=0$ can be evaluated directly.

Every query has polynomial length. The gadget has $O(p^2)$ incidences, while the sum construction has polynomial size in the original formula length and $\log p$. In the hypersurface construction the least power satisfies $Q\le p(u+2)$, so its explicitly listed $Q+1$ variables and its arithmetic circuit still have polynomial size because $p\le M$. Thus the reduction also supplies valid polynomial-size circuit inputs to the exact-height problem.

Finally, any $\#\mathrm P$ function counts the accepted strings of a polynomial-time verifier of polynomially bounded length. Unrolling that verifier into a Boolean circuit and applying the unique-extension gate clauses just used gives a parsimonious width-three formula. The reduction consequently applies to all $\#\mathrm P$ functions. Replacing each decision query by an exact-height query and testing whether the answer is two proves the final assertion. The value $\infty$ is simply one of the negative answers. The finite-height bound gives $O(d)$ bits for every finite answer, so reading a height answer also costs polynomial time.
\end{proof}
The prime varies in this theorem. It supplies no $\#\mathrm P$-hardness statement at one fixed characteristic, and its exact-counting equivalence concerns the uniform height-two oracle. The complexity of the unrestricted full-height function still requires an upper bound matching its lower bound.


\subsection{Height supplied in unary}
The matrix criterion also permits modular counting when the requested
height grows with the input. The relevant witness is a path in the
matrix, rather than a coefficient of degree exponential in that height.

\begin{theorem}[Unary height decisions]\label{thm:unary-height}
Fix a prime $p$. On the smooth homogeneous circuit inputs specified
above, with a positive integer $r$ supplied in unary, deciding
$h(X)=r$ is $\mathrm{Mod}_p\mathrm P$-complete under polynomial-time
many-one reductions. In particular, membership holds uniformly in $r$;
the completeness assertion for this variable-$r$ problem already holds
on its restriction $r=2$.

Let $H_{\mathrm{unary}}$ denote the same decision problem with both
$p$ and $r$ supplied in unary, and with $p$ promised prime. Then
$H_{\mathrm{unary}}\in\mathrm P^{\#\mathrm P}$, using one exact
counting query followed by a remainder computation. With
promise-respecting oracle calls,
\[
 \mathrm{FP}^{H_{\mathrm{unary}}}=\mathrm{FP}^{\#\mathrm P}.
\]
\end{theorem}
\begin{proof}
We first construct nonnegative counts for the coefficients entering
the operator in \eqref{eq:quasi-delta}. Homogenize the input circuit
through degree $d$ and keep its degree-$d$ output. Evaluate each
degree-zero subcircuit in $\mathbb F_p$ and replace it by its residue
in $\{0,\ldots,p-1\}$. Rewrite subtraction as addition after
multiplication by $p-1$. The resulting arithmetic circuit uses only
addition, multiplication and nonnegative integer scalar literals;
it represents an integral homogeneous lift $\widetilde F$ of $F$.
Repeated squaring constructs a circuit for
\[
 G=\widetilde F^{p-1}=\sum_{|\beta|=D}c_\beta x^\beta,
 \qquad D=d(p-1),\qquad c_\beta\ge0.
\]
Its size is polynomial in the input length, and $D$ is polynomial
when $p$ is fixed or unary.

A parse witness at an addition chooses one child, at a product takes
both children, and at a scalar $b$ chooses one of $b$ labels. The
exponent vector records the variable leaves. Thus the number of
witnesses with vector $\beta$ is exactly $c_\beta$. These are
polynomial-length witnesses even though the circuit may share
subcomputations. After degree-zero gates have been collapsed, every
positive-degree leaf contributes at least one to total degree. There
are at most $D$ such leaves and at most $D-1$ product nodes having
two positive-degree children. Between those branching nodes a chain
has length at most the circuit size; scalar children add at most
one leaf per node. Hence a parse has polynomially many nodes,
bounded by the circuit size times $O(D)$. Choose a fixed-length
encoding with unique padding, so each parse contributes exactly one
accepting witness and the witnesses have a canonical lexical order.
This proves that $c_\beta$ is a $\#\mathrm P$ count, uniformly in
the indicated input parameters.

Consider the integral polynomial
\[
 Q=\frac{G^p-\sum_{|\beta|=D}c_\beta^{\,p}x^{p\beta}}p.
\]
The sum is over distinct collected monomials. For an exponent vector
$\gamma$, the numerator coefficient counts ordered $p$-tuples of
parse witnesses whose vectors sum to $\gamma$, excluding tuples
whose exponent vectors are all equal. The excluded number is
precisely $\sum_{p\beta=\gamma}c_\beta^p$: it includes all parses
of the same monomial, not only tuples of identical parses.

Cyclic rotation acts freely on the remaining tuples. Indeed, since
$p$ is prime, a tuple fixed by a nonidentity rotation has all its
parse witnesses equal, and therefore has all exponent vectors equal.
Consequently $[x^\gamma]Q$ is the number of these free orbits.
A verifier guesses the $p$ parses, checks their exponent sum and
that their vectors are not all equal, and accepts exactly the
lexically least cyclic rotation. Both the witness and the comparison
with all $p$ rotations have polynomial length and running time for
fixed or unary $p$. Every coefficient of $Q$ is thus itself a
$\#\mathrm P$ count. As a small example, if $G=cx^\beta+ex^\eta$
with $\beta\ne\eta$ and $p=2$, then
$Q=cex^{\beta+\eta}$: the two possible orders of one parse of
each monomial form one orbit.

By coefficient-lift independence in \eqref{eq:quasi-delta},
$Q\bmod p=\Delta$. In particular a coefficient of $Q$ gives a
nonnegative counting lift of each matrix entry
\[
 M_{\alpha\beta}=[x^{p\alpha+a-\beta}]\Delta.
\]
A negative component of $p\alpha+a-\beta$ gives zero witnesses.
For $j\ge1$, expand the scalar
\[
 y_j=\ell(M^{j-1}g)
\]
as a sum over paths of $j-1$ matrix transitions ending at the
monomial $x^a$. A path witness specifies all its degree-$D$
monomial indices, a parse witnessing its initial coefficient of
$G$, and an orbit witness for each transition. Each index has
$O(d\log(D+1))$ bits. Therefore, for unary $j$, this is a
polynomial-length witness, with a polynomial-time verifier.
Let $N_j$ be its number of accepting witnesses. Summing over paths
and multiplying the choices along each path proves
\[
 N_j\equiv y_j\pmod p.
\]
This construction includes $j=1$, when the path has no transitions
and the count is the coefficient of $x^a$ in $G$.

Nonnegative counting functions are closed under sums and under
polynomially many products: a witness selects a summand or supplies
one witness for each factor, respectively. Thus
\begin{equation}\label{eq:unary-count}
 Z_r=N_r^{p-1}\prod_{j=1}^{r-1}
       \bigl(1+(p-1)N_j^{p-1}\bigr)
\end{equation}
is a $\#\mathrm P$ count when $p$ is fixed or unary and $r$ is
unary. The sum of the path lengths in this product is $O(r^2)$,
multiplied by polynomial factors in the unary prime and circuit
length. Fermat's identity in $\mathbb F_p$ gives
\[
 Z_r\not\equiv0\pmod p
 \quad\Longleftrightarrow\quad
 y_1=\cdots=y_{r-1}=0\ \hbox{ and }\ y_r\ne0.
\]
The matrix criterion identifies this condition with $h(X)=r$.
For fixed $p$ this proves $\mathrm{Mod}_p\mathrm P$ membership;
the height-two reduction proves completeness of the variable-$r$
problem. If both $p$ and $r$ are unary, an exact counting oracle
computes $Z_r$, whose binary length is polynomial in the witness
bound, and the remainder modulo $p$ decides the problem. This
proves the claimed $\mathrm P^{\#\mathrm P}$ upper bound.
Conversely, every query to the promised uniform height-two oracle
in Theorem~\ref{thm:uniform-counting} is a valid query to
$H_{\mathrm{unary}}$ with $r=2$. That theorem recovers every
$\#\mathrm P$ function. Composing these two simulations proves
the stated equality of function-oracle classes.
\end{proof}

\begin{corollary}[Unary thresholds]\label{cor:unary-threshold}
For a fixed prime $p$ and unary $R$, deciding $h(X)\le R$ belongs
to $\mathrm{Mod}_p\mathrm P$. If a specified polynomial $q$ is
fixed and the input additionally promises that every finite height
is at most $q(L)$, where $L$ is circuit input length, then the
exact-height function, including $\infty$, belongs to
$\mathrm{FP}^{\mathrm{Mod}_p\mathrm P}$ on that promise.
\end{corollary}
\begin{proof}
Over $\mathbb F_p$, the expression
$1-\prod_{j=1}^R(1-y_j^{p-1})$ is nonzero exactly when some
$y_j$ is nonzero. Replace each negative scalar by its residue
$p-1$ and use the counts just constructed; the result is a
$\#\mathrm P$ count modulo $p$ with polynomial witness length.
For the second assertion, query the exact-height predicate for
$r=1,\ldots,q(L)$. Return its first positive answer or $\infty$
if none occurs. The additional promise makes the latter conclusion
valid. No polynomial bound on finite heights is assumed for
general circuit inputs, and the argument does not allow a
binary-encoded height parameter.
\end{proof}


\paragraph{What the lower bounds establish.}
The counting construction makes height one
$\mathrm{Mod}_p\mathrm P$-complete for every fixed prime.  Combined
with the fixed-gadget reduction above, the preceding
membership theorem makes height two $\mathrm{Mod}_p\mathrm P$-complete
as well.  At each fixed characteristic, the full exact-height function is therefore
$\mathrm{Mod}_p\mathrm P$-hard. With the characteristic as part of the
input, Theorem~\ref{thm:uniform-counting} strengthens the lower bound
to $\#\mathrm P$-hardness under nonadaptive polynomial-time oracle
reductions. The full function lies in $\mathrm{FPSPACE}$, and hence in
$\mathrm{FEXPTIME}$, the function class with deterministic running time
$2^{\operatorname{poly}(L)}$ for input length $L$.
These are distinct bounds; they do not constitute a complete complexity
classification of the full function. A single height-one zero test
does not determine the full residue of a counting function.
Theorem~\ref{thm:unary-height} supplies uniform modular counting for
unary height parameters, but does not decide infinite height in that
class or extend the witness bound to binary height parameters.
